\documentclass{article}
\usepackage{graphicx} % Required for inserting images
\usepackage{tcolorbox}
\usepackage{amsthm}
\usepackage{amsmath}
\usepackage{amssymb}
\usepackage{enumitem}
\usepackage{hyperref}
\usepackage{braket}

\title{Quantum Computation Notes \\ \large Hidden Subgroup Problem}
\author{}
\date{}
\begin{document}
\maketitle
\section{Preliminaries}
Consider abelian group $G$. Then, a \textbf{representation} $\rho$ of $G$ can be defined as a function that maps $G$ to a matrix group, preserving group multiplication. The \textbf{character} of a matrix group $G \subset M_n$ is a function defined by $\chi(g) = tr(g)$ for $g \in G$. More broudly, a function is a character if it is a mapping from $G$ to nonzero complex numbers that is a group homomorphism: i.e $f(g_1 g_2) = f(g_1) f(g_2)$. The characters of an Abelian group are linear.\\
The \textbf{character group} or \textbf{dual group} is the group under pointwise multiplication consisting of all the characters of group $G$ and is denoted $\hat{G}$. For example, consider group $\mathbb{Z}_7$. This group is generated by 1, so the behavior of the character is determined by $\chi(1)$. Then, we see it must be that $\chi(1)^7 = \chi(1)$ since in $\mathbb{Z}_7, 1 \times 7 = 1$. We can see that:
$$\chi_k(n) = e^{2\pi i kn/7}$$
satisfies the definition of a character group for $\mathbb{Z}_7$.\\
If you define inner product on functions $f,g$:
$$\braket{f|g} = \frac{1}{|G|} \sum_{x \in \hat{G}} f(x) \overline{f(x)}$$
We then define orthogonality:
$$\braket{\chi_x(a)|\chi_x(b)} = \frac{1}{|G|} \sum_{x \in \hat{G}} \chi_x(a) \overline{\chi_x(b)} = \delta_{ab}$$
equivalently:
$$\frac{1}{|G|} \sum_{x \in G} \chi_y(x) \overline{\chi_{y'} (x)} = \delta_{y, y'}$$
\section{Quantum Fourier Transform}
Consider abelian group $G$. Then, the Fourier transform on $G$ is defined by:
$$F_G := \frac{1}{\sqrt{G}} \sum_{x \in G} \sum_{y \in \hat{G}} \chi_y(x) \ket{y}\bra{x}$$
Since $G, \hat{G}$ are isomorphic, we label elements of $\hat{G}$ with elements of $G$. 
\begin{tcolorbox}\textbf{The Fourier Transform on G is unitary.}
	\begin{proof}
		We show that $\braket{u|F^\dag_g F_g|v} = \delta_{uv}$.
		We see that:
		\begin{align*}
			F_g\ket{v} &= \frac{1}{\sqrt{|G|}} \sum_{x \in G} \sum_{y \in \hat{G}} \chi_y(x)\ket{y} \braket{x|v}\\
				   &= \frac{1}{\sqrt{|G|}} \sum_{y \in \hat{G}} \chi_y(v) \ket{y}
		\end{align*}
		and that:
		\begin{align*}
			\bra{u}F^\dag_g &= \frac{1}{\sqrt{|G|}} \sum_{a \in G} \sum_{b \in \hat{G}} \overline{\chi_b(a)} \braket{u|a}\bra{b}\\
					&= \frac{1}{\sqrt{|G|}} \sum_{b \in \hat{G}} \overline{\chi_b(u)} \bra{b}
		\end{align*}
		Then:
		\begin{align*}
			\braket{u|F^\dag_g F_g|v} &= \frac{1}{|G|} \sum_{y\in\hat{G}}\sum_{b \in \hat{G}}\chi_y(v) \overline{\chi_b(u)} \braket{b|y}\\
						  &= \frac{1}{|G|} \sum_{y \in \hat{G}} \chi_y(v)\overline{\chi_y (u)} = \delta_{uv}
	\end{align*}
	\end{proof}
\end{tcolorbox}
The construction of the circuit for QFT for $\mathbb{Z}_N$ is identical to in Nielsen and Chuang. The action of the QFT can be expressed in product-state form and this clears up the construction of the circuit. 
The QFT over $\mathbb{Z}_{2^n}$ can be used for phase estimation. Phase estimation, given some unitary with eigenvector $\ket{u}$ and eigenvalue $e^{2 \pi i \phi}$
$$\ket{u,0} \xrightarrow[]{H} \sum_{x \in \mathbb{Z}_{2^n}}\ket{u, x} \xrightarrow[]{C(U)} \sum_{x} e^{2\pi i \phi} \ket{u,x} \xrightarrow[]{iQFT} \ket{u, \hat{\phi}}$$
\section{QFT over Finite Abelian Groups}
Suppose we implement QFT over arbitrary cyclic group $\mathbb{Z}_N$:
$$F_{\mathbb{Z}_N} = \frac{1}{\sqrt{N}} \sum_{x,y \in \mathbb{Z}_N} \omega^{xy}_N \ket{y}\bra{x}$$
The circuit we constructed in Nielsen and Chuang using $R_{k}$ circuits (as far as I recall) only works for $N$ a power of two. We can realize $F_{\mathbb{Z}_{N}}$ using phase estimation. We want to perform transformation that maps $\ket{x} \to \ket{\hat{x}}$ where $\ket{\hat{x}} := F_{\mathbb{Z}_N} \ket{x}$ denotes a Fourier basis state. It is easy to perform $\ket{x,0} \to \ket{x, \hat{x}}$ using some unitary that performs $U\ket{x, 0} \to \ket{x, f(x)}$. Then, consider:
$$ U := \sum_{x \in \mathbb{Z}_N} \ket{x+1}\bra{x}$$
The eigenstates are $\ket{\hat{x}} = F_{\mathbb{Z}_N} \ket{x}$ since we can see that:
$$F^{\dag}_{\mathbb{Z}_N} U F_{\mathbb{Z}_N} = \sum_{x \in \mathbb{Z}_N} \omega_N^{-x} \ket{x}\bra{x}$$
Then, using phase estimation on $U$, we can perform:
$$\ket{\hat{x},0}\to\ket{\hat{x},x}$$
If we perform the transpose of $QPE$, we can eliminate $x$. Thus:
$$\ket{0,x} \xrightarrow[]{} \ket{\hat{x}, x} \xrightarrow[]{QPE^\dag} \ket{\hat{x}, 0}$$
\section{Discrete Log}
Let $G = \braket{g}$ by a cyclic group generated by $g$ multiplicatively. Given element $x \in G$ the discrete logarithm of $x \in G$ with respect to $g$ is $\log_g x$, the smallest non-negative integer $g^\alpha = x$. For example, ECDLP, or $G = \mathbb{Z}_7^X, \log_3 2 = 2$ since $3^2 = 9 \bmod 7 = 2$.
The hardness of discrete log is used for Diffie-Hellman exchange. Suppose we work in $\mathbb{Z}_p^x$. Then the protocol works like this:
\begin{enumerate}
	\item Alice and Bob choose $p$ (prime) and high order integer $g$ and for the sake of simplicity let $g$ be a generator for $\mathbb{Z}_p^X$.
	\item Alice chooses $a \in Z_{p-1}$ and computes $A = g^a \bmod p$.
	\item Bob chooses $b \in Z_{p-1}$ and computes $B=g^b \bmod p$.
	\item Alice computes $K = B^a \bmod p = g^ab \bmod p$.
	\item Then, both can compute $K = B^a \bmod p$ or $K = A^b \bmod p$.
	They can show $K$ and an onlooker only sees $p,g, A, B$.
\end{enumerate}
\section{Hidden Subgroup Problem}
Hidden subgroups can be efficiently calculated on a quantum computer. Discrete log, factoring, etc. are all examples of hidden subgroup problems. We are given black box function $f: G \to S$ where $G$ is a known group and $S$ is a finite set. The function satisfies:
$$f(x) = f(y) \iff x^{-1} y \in H \iff y = xh, h \in H$$
for unknown subgroup $H \leq G$. W want to find the generators of $H$. In other words, $f$ is constant on cosets of $H$ and distinct on different cosets. A simple example like $\mathbb{Z}_6$ makes it clear.\\
\textbf{Theorem 5.1} Suppose $G$ has a set $\mathcal{H}$ of $N$ subgroups whose only common element is the identity. Then, to solve the HSP, a deterministic classical computer must make $\Omega(\sqrt{N})$ queries.
\begin{proof} 
	Suppose the oracle behaves adversarially as follows:\\
	On the $l$th query, the algorithm queries $g_l$ which we assume is different from $g_1, \dots, g_{l-1}$. If there is some subgroup $H \in \mathcal{H}$ such that $g_k \not \in g_j H, 1 \leq j \leq k \leq l$, then the oracle outputs $l$, or else the oracle outputs a generating set for $H$. Basically iwfor all the group elements we have checked so far, we haven't found some points in constant coset. Suppose there is an algorithm that queries the oracle $t$ times before forcing the oracle to give up. Up to $t$, there are at most $t(t-1) $ values of $g_k g_j^{-1}$ for $N$ possible subgroups. To satisfy $N$ for all $H\in \mathcal{H}$, there is some pair $j,k$ such that $g_k g_j^{-1} \in H$ we must check all the $t(t-1) \geq N \implies t = \Omega(\sqrt{N})$.
\end{proof}
\section{Shor's Algorithm for DLP}
Shor's algorithm can be used to calculate discrete logarithms in $\mathbb{Z}_N^\times.$ Suppose we are given $x \in G$ where $G$ is a cyclic group generated by $g$. We want to calculate $x = g^\alpha$. Suppose we know the order of $G$, $N = |G|$. We can learn this using order-finding. \\
The discrete log is abelian HSP in group $\mathbb{Z}_N \times \mathbb{Z}_N$. Define:
$$f(a,b) = x^a g^b$$
Then, since $f(a,b) = g^{\alpha a + b}$, $f$ is constant on the lattice defined by:
$$\mathcal{L}_\lambda := \{(\alpha, \beta) \in \mathbb{Z}_N^2 : \alpha a + b = \lambda\}$$
It is a line. Equivalently, the subgroup is:
$$H = L_0 = \{ (0,0), (1, -\alpha), (2, -2\alpha), \dots\}$$
We can naturally see the tuple period interpretation of $f(a,b) = g^{\alpha a + b} = g^{\alpha(a-l) + (b+\alpha l)} = f(a-l, \alpha l + b)$ so we have period tuple $(-l, \alpha l)$
and $\mathcal{L}_\lambda$ is the $\lambda$ coset, or $\lambda + H$.\\
More precisely, the cosets of $H$ in $\mathbb{Z}_N \times \mathbb{Z}_N$ are of the form $(\gamma, \delta) + H$ with $\gamma, \delta \in \mathbb{Z}_N$. The set of cosets defined for each $\delta \in \mathbb{Z}_N$ form a complete set of cosets:
$$(0, \delta) + H = \{ (\alpha, \delta - \alpha \log_g(x)) | \alpha \in \mathbb{Z}_N \}$$
The algorithm roughly proceeds:
\begin{enumerate}
	\item Place in even superposition $$\frac{1}{N} \sum_{\alpha, \beta \in \mathbb{Z}_N} \ket{\alpha, \beta, 0}$$
	\item Apply controlled $f$ gate on third register and measure.
	$$\frac{1}{N} \sum{\alpha, \beta}\ket{\alpha, \beta, f(\alpha, \beta)}$$
	When we measure, it collapses to some coset. $\ket{(0,\delta) + H} = \ket{L_\delta} = $
	$$\frac{1}{\sqrt{N}} \sum_{\alpha \in \mathbb{Z}_N} \ket{\alpha, \delta-\alpha \log_g x}$$
\item We now can perform a QFT on the two registers:
	$$\frac{1}{N^{3/2}} \sum_{\alpha, u, v \in \mathbb{Z}_N} \omega_N^{u\alpha + v(\delta - \alpha \log_g x)} \ket{u,v}= \frac{1}{N^{3/2}} \sum_{u,v} \omega_N^{uv} \sum_{\alpha \in \mathbb{Z}_N} \omega_N^{\alpha(u-v \log_g x)} \ket{u,v}$$
	Using $\sum_{\alpha \in \mathbb{Z}_N} \omega^{\alpha \beta}_N = N \delta_{\beta 0}$ (easy to see using Taylor series):
	$$ = \frac{1}{\sqrt{N}} \sum_{v \in \mathbb{Z}_N} \omega^{v \delta}_N \ket{v \log_g x, v}$$
\end{enumerate}
\section{The Abelian HSP}
We describe the Hidden Subgroup over some arbitrary finite abelian group. We will use additive notation, so $f$ hide $H$ as $f(x) = f(y) \iff x - y \in H$. The algorithm generally follows Shor's algorithm. 
\begin{enumerate}
	\item Create superposition:
		$$\ket{G} = \frac{1}{\sqrt{|G|}} \sum_{x \in G} \ket{x}$$
	\item Compute the function in another register.
		$$\xrightarrow[]{C(U)} \frac{1}{\sqrt{|G|}}\sum_{x \in G}\ket{x, f(x)}$$
	\item Measure the second register, collapsing the first to a uniform superposition of a coset.
		$$ \ket{x+H} = \frac{1}{\sqrt{|H|}} \sum_{h \in H} \ket{x+h}$$
	\item Perform a Fourier transform over G. We see that:
		$$F_g \ket{x+H} = \hat{\ket{x+H}}  = \frac{1}{\sqrt{|H||G|}} \sum_{y \in \hat{G}} \sum_{h \in H} \chi_y(x + h)\ket{y}$$
	If we define $\chi_y(H) := \frac{1}{|H|} \sum_{h \in H} \chi_y(h)$,
	$$= \sqrt{\frac{|H|}{|G|}} \sum_{y \in \hat{G}} \chi_y(x) \chi_y(H) \ket{y}$$
	We can see two things: firstly, that if $\chi_y(h) = 1$ for all $h \in H$, then clearly $\chi_y(H) = 1$. Secondly if there is $h' \in H$ such that $\chi_y(h') \neq 1$ then since $h' + H = H$ (group closed)
	\begin{align*}
		\chi_y(H) &= \frac{1}{|H|} \sum_{h \in h' + H} \chi_y(h)\\
			  &= \frac{1}{|H|} \sum_{h \in H} \chi_y(h' + h)
	\intertext{Since characters are group homomorphisms:}
			  &= \frac{1}{|H|} \sum_{h \in H} \chi_y(h') \chi_y(h)\\
			  &= \chi_y(h') \chi_y (H)\\
			  \intertext{Then, since $\chi_y(H) = \chi_y(h') \chi_y(H)$ for nonzero $\chi_y(h')$:}
			  &\implies \chi_y(H) = 0
\end{align*}
This alternatively follows from orthogonality of characters.
Therefore only the terms of $y \in \hat{G}$ where all $\chi_y(h) = 1$ for all $h \in H$ remain.
$$F_G\ket{x+H} = \sqrt{\frac{|H|}{|G|}} \sum_{y: \chi_y (H) = 1} \chi_y(x) \ket{y}$$
\item Measure the register and receive $y_0$. We then can determine some character $\chi_{y_0}$ where for all $h \in H$, $\chi_{y_0}(h) = 1$. 
	$$ H \subset \text{ker}(\chi_{y_0}) = \{g \in G | \chi_{y_0}(g) = 1\}$$
	By measuring several $y_0, \dots, y_k$ we can take the intersection of the various kernels and very quickly determine $H$. We show we only need $O(\log|G|)$ measurements to determine $H$ with high accuracy.
	\begin{proof}
		Suppose that at some point the intersection of kernels is not $H$:
		$$H \neq \text{ker} \chi_{y_0} \cap \dots \cap \text{ker} \chi_{y_k} = K \leq G$$
		where $K$ is a subgroup of $G$ since the intersection of subgroups is a subgroup. We also see $H< K$. Then, we see that the probability of sampling $\chi_y$ s.t. $\chi_y(H) = 1$ has probability $|H|/|G|$, originating from the normalization factor of $\sqrt{|H|/|G|}$ in the state of superpositions of $\chi_y(H) = 1$. Then, we want the probability that we sample $y$ such that $\chi_y$ satisfies $K \leq \text{ker} \chi_y$ i.e. that we do not learn anything from the new character. 
		$$\frac{|H|}{|G|} |\{ y\in \hat{G} | K \leq \text{ker} \chi_y\}|$$
		The number of such y's is $|G|/|K|$, since if subgroup $K$ was hidden, we could sample such $y$'s uniformly with probability $|K|/|G|$. Therefore the probability is:
		$$\frac{|H|}{|K|} \leq \frac12$$
		where the inequality comes from $|K| \leq 2|H|$ by Lagrange's theorem. If we sample $y$ where $K \not \leq \text{ker} \chi_y$, then the new $K$ (that is, intersected with the kernel of the new character) $|K \cap \text{ker} \chi_y| \leq |K|/2$ (which happens with probability greater than $\frac12$). Thus if we repeat the process $O(\log |G|)$ times, it is extremely likely the resulting subgroup is $H$.
	\end{proof}
\end{enumerate} 
\end{document}
